\section{Exact path-completion calculation}
\label{sec:path-completion}

A second standard attempt hides the endpoint of a long path.  The exact
Green function again shows that its distinguishability depth never exceeds
the output scale.

Let $P_N$ be a path of $N$ edges with the seed at one endpoint, and put
\begin{equation}
 r:=\frac{1-\sqrt\alpha}{1+\sqrt\alpha},
 \qquad
 t:=-\log r
 =2\operatorname{artanh}(\sqrt\alpha).
 \label{eq:path-decay}
\end{equation}

\begin{proposition}[Path completion sensitivity; \ProvedStatus]
\label{prop:path-completion}
The normalized seed value on $P_N$ is
\begin{equation}
 y_v^0(P_N)
 =\sqrt\alpha\,
  \frac{1+r^{2N}}{1-r^{2N}}
 =\sqrt\alpha\,\coth(Nt).
 \label{eq:path-seed-value}
\end{equation}
Consequently
\begin{equation}
 y_v^0(P_N)-y_v^0(P_{2N})
 =\frac{\sqrt\alpha}{\sinh(2Nt)}.
 \label{eq:path-completion-gap}
\end{equation}
If the two rooted paths must have seed values separated by more than
$2\epsppr$, then necessarily
\begin{equation}
 N
 <\frac{1}{2t}
   \operatorname{arsinh}\!\left(
      \frac{\sqrt\alpha}{2\epsppr}
   \right)
 \leq\frac1{8\epsppr}.
 \label{eq:path-depth-output-scale}
\end{equation}
Thus this indistinguishable-completion pair cannot prove more than
$O(1/\epsppr)$ queries.
\end{proposition}

\begin{proof}
The nonsource recurrence associated with \cref{eq:killed-system} has roots
$r$ and $r^{-1}$; the reflecting degree-one endpoint contributes the image
term $r^{2N-k}$.  Enforcing the sourced endpoint row gives
\cref{eq:path-seed-value}.  The identity
$\coth z-\coth(2z)=1/\sinh(2z)$ gives
\cref{eq:path-completion-gap}.  Separation by more than $2\epsppr$ implies
the first inequality in \cref{eq:path-depth-output-scale}.  Finally
$t=2\operatorname{artanh}(\sqrt\alpha)\geq2\sqrt\alpha$ and
$\operatorname{arsinh}(u)\leq u$ for $u\geq0$.
\end{proof}

The precise regimes are informative.  When
$N\sqrt\alpha\ll1$, the gap is $\Theta(1/N)$.  At
$N=\Theta(1/\sqrt\alpha)$ it is $\Theta(\sqrt\alpha)$, and beyond that scale
it decays as
$\sqrt\alpha\exp(-\Theta(N\sqrt\alpha))$.  If
$\epsppr\ll\sqrt\alpha$, one may hide the endpoint for an additional
logarithmic factor, but
$\alpha^{-1/2}\log(\sqrt\alpha/\epsppr)=O(1/\epsppr)$.
